\documentclass[xetex,mathserif,serif]{beamer}

\usepackage{xunicode}
\usepackage{xltxtra}
\usepackage{color}
\usepackage{url}
\usepackage{listings}
\usepackage{fontspec}
\usepackage{geometry}
\usepackage{lastpage}
\usepackage{fancyhdr}
\usepackage{amsmath}
\usepackage{amsthm}
\usepackage{amssymb}
\usepackage{blkarray}
\usepackage{multicol}
\usepackage{relsize}

\definecolor{solarized@base03}{HTML}{002B36}
\definecolor{solarized@base02}{HTML}{073642}
\definecolor{solarized@base01}{HTML}{586e75}
\definecolor{solarized@base00}{HTML}{657b83}
\definecolor{solarized@base0}{HTML}{839496}
\definecolor{solarized@base1}{HTML}{93a1a1}
\definecolor{solarized@base2}{HTML}{EEE8D5}
\definecolor{solarized@base3}{HTML}{FDF6E3}
\definecolor{solarized@yellow}{HTML}{B58900}
\definecolor{solarized@orange}{HTML}{CB4B16}
\definecolor{solarized@red}{HTML}{DC322F}
\definecolor{solarized@magenta}{HTML}{D33682}
\definecolor{solarized@violet}{HTML}{6C71C4}
\definecolor{solarized@blue}{HTML}{268BD2}
\definecolor{solarized@cyan}{HTML}{2AA198}
\definecolor{solarized@green}{HTML}{859900}
\definecolor{yaleblue}{HTML}{0E4C92}

\setbeamertemplate{navigation symbols}{}
% \setbeamerfont{title}{family=\old}
% \setbeamerfont{author}{family=\tfont}%
% \setbeamerfont{frametitle}{family=\oldA}
% \setbeamerfont{date}{family=\dfont}

\setbeamertemplate{itemize items}{--}
\setbeamercolor*{item}{fg=black}

\defaultfontfeatures{Mapping=tex-text}
\hypersetup{pdfstartview={FitH}}

\newcommand{\old}[1]{\fontspec[Alternate=1,Ligatures={Common}]{Hoefler Text}\fontsize{18pt}{30pt}\selectfont #1}%
\newcommand{\oldA}[1]{\fontspec[Alternate=1,Ligatures={Common, Rare}]{Hoefler Text}\fontsize{12pt}{15pt}\selectfont #1}%
\newcommand{\oldB}[1]{\fontspec[Ligatures={Common}]{Didot}\fontsize{12pt}{15pt}\color{solarized@base02}\selectfont #1}%
\newcommand{\tfont}[1]{\fontspec[Alternate=1,Ligatures={Common}]{Hoefler Text}\fontsize{12pt}{20pt}\selectfont #1}%
\newcommand{\dfont}[1]{\fontspec[Ligatures={Common}]{Didot}\fontsize{12pt}{12pt}\selectfont #1}%

\newcommand{\minimize}{\mathop{\mathrm{minimize}}}
\newcommand{\argmin}{\mathop{\mathrm{arg\,min}}}
\newcommand{\argmax}{\mathop{\mathrm{arg\,max}}}
\newcommand{\st}{\mathop{\mathrm{subject\,\,to}}}

\newcommand\independent{\protect\mathpalette{\protect\independenT}{\perp}}
\def\independenT#1#2{\mathrel{\rlap{$#1#2$}\mkern2mu{#1#2}}}

\setlength{\parindent}{0pt}
\setlength{\parskip}{12pt}

\setromanfont [Ligatures={Common}, Numbers={OldStyle}, Variant=01,
 BoldFont={LinLibertine_RB.otf},
 ItalicFont={LinLibertine_RI.otf},
 BoldItalicFont={LinLibertine_RBI.otf}
 ]{LinLibertine_R.otf}



\begin{document}

%%%%%%%%%%%%%%%%%%%%%%%%%%%%%%%%%%%%%%%%%%%%%%%%%%%
\begin{frame}[fragile] \frametitle{}

\vfill

{\fontsize{0.7cm}{0cm}\selectfont Lecture 07 \\\vspace{0.2cm}
Hypothesis Testing with Multivariate Regression}\\\vspace{0.5cm}
23 September 2015

\vspace{2cm}

\begin{minipage}{0.6\textwidth}
Taylor B. Arnold \\
Yale Statistics \\
STAT 312/612
\end{minipage}
\hfill
\begin{minipage}{0.3\textwidth}\raggedleft
\includegraphics[scale=0.3]{../yale-logo.png}
\end{minipage}%

\end{frame}

%%%%%%%%%%%%%%%%%%%%%%%%%%%%%%%%%%%%%%%%%%%%%%%%%%%
\begin{frame}[fragile] \frametitle{}

{\color{yaleblue}\fontsize{16pt}{20pt}\selectfont Goals for today}

\begin{enumerate}
\item Galton heights data
\item Hypothesis tests
\item F and T simulations
\end{enumerate}

\end{frame}

%%%%%%%%%%%%%%%%%%%%%%%%%%%%%%%%%%%%%%%%%%%%%%%%%%%
\begin{frame}[fragile] \frametitle{}

\begin{flushright}
{\color{yaleblue}\sc\fontsize{1cm}{0cm}\selectfont Linear models assumptions}
\end{flushright}

\end{frame}

%%%%%%%%%%%%%%%%%%%%%%%%%%%%%%%%%%%%%%%%%%%%%%%%%%%
\begin{frame}[fragile] \frametitle{}

{\bf Hypothesis tests}

Consider testing the hypothesis $H_0: beta_j = b$ for
some particular $b$ and $j$.

\pause Under assumptions I-V we have the following:
\begin{align*}
\left. \widehat{\beta} - b \right| X, H_0 \sim \mathcal{N} ( 0, \sigma^2 \left( (X^t X)^{-1}_{jj} \right))
\end{align*}

\end{frame}


%%%%%%%%%%%%%%%%%%%%%%%%%%%%%%%%%%%%%%%%%%%%%%%%%%%
\begin{frame}[fragile] \frametitle{}

{\bf Z-test}

This suggests the following test statistic:
\begin{align*}
z &= \frac{\widehat{\beta} - b}{\sqrt{\sigma^2  \left( (X^t X)^{-1}_{jj} \right)}}
\end{align*}
With,
\begin{align*}
\left. z \right| X, H_0 \sim \mathcal{N} (0, 1)
\end{align*}

\end{frame}

%%%%%%%%%%%%%%%%%%%%%%%%%%%%%%%%%%%%%%%%%%%%%%%%%%%
\begin{frame}[fragile] \frametitle{}

{\bf T-test}

As in the simple linear linear regression case, we need to estimate
$\sigma^2$ with $s^2$. This yields the following test statistic:
\begin{align*}
t &= \frac{\widehat{\beta} - b}{\sqrt{s^2  \left( (X^t X)^{-1}_{jj} \right)}} \\
&= \frac{\widehat{\beta} - b}{\text{S.E.}(\widehat{\beta}_j)}
\end{align*}


\end{frame}

%%%%%%%%%%%%%%%%%%%%%%%%%%%%%%%%%%%%%%%%%%%%%%%%%%%
\begin{frame}[fragile] \frametitle{}

{\bf T-test}

The test statistic has a $T$-distribution  with $(n-p)$ degrees
of freedom
\begin{align*}
\left. t \right| X, H_0 \sim t_{n-p}
\end{align*}
\pause This time around, we'll actually prove this.

\end{frame}

%%%%%%%%%%%%%%%%%%%%%%%%%%%%%%%%%%%%%%%%%%%%%%%%%%%
\begin{frame}[fragile] \frametitle{}

Note that:
\begin{align*}
t &= \frac{\widehat{\beta} - b}{\sqrt{\sigma^2  \left( (X^t X)^{-1}_{jj} \right)}} \cdot \sqrt{\frac{\sigma^2}{s^2}} \\
&= \frac{z}{\sqrt{s^2 / \sigma^2}} \\
&= \frac{z}{q / (n-p)}
\end{align*}
\pause Where $q = r^t r / \sigma^2$.

\pause We need to show that (1) $q | X \sim \chi^2_{n-p}$ and
(2) $z \independent q | X$.

\end{frame}

%%%%%%%%%%%%%%%%%%%%%%%%%%%%%%%%%%%%%%%%%%%%%%%%%%%
\begin{frame}[fragile] \frametitle{}

{\bf Lemma} If $B$ is an idempotent matrix and
$v \sim \mathcal{N} (0, \mathbb{I}_n)$, then
$v^t B v \sim \chi^2_{\text{rank(B)}}$.

\pause
See \url{http://www2.econ.iastate.edu/classes/econ671/hallam/documents/QUAD_NORM.pdf}
for a proof of this lemma.

\end{frame}

%%%%%%%%%%%%%%%%%%%%%%%%%%%%%%%%%%%%%%%%%%%%%%%%%%%
\begin{frame}[fragile] \frametitle{}

(1) We know that $r^t r = \epsilon^t M \epsilon$, so:
\begin{align*}
q &= \frac{r^t r}{\sigma^2} \\
&= \frac{\epsilon^t}{\sigma} \cdot M \cdot \frac{\epsilon}{\sigma}
\end{align*}
\pause From the lemma then $q \sim \chi^2_{\text{rank(M)}}$. Finally,
noting that $rank(M) = tr(M) = n - p$ finishes the first part of the
proof.

\end{frame}

%%%%%%%%%%%%%%%%%%%%%%%%%%%%%%%%%%%%%%%%%%%%%%%%%%%
\begin{frame}[fragile] \frametitle{}

(2) The random variables $\widehat{\beta}$ and $r$ are both
linear combinations of $\epsilon$, and are therefore jouintly
normally distributed. As they are uncorrelated (problem set 2),
this implies that they are independent. Finally, this implies that
$z = f(\widehat{\beta})$ and $t = g(r)$ and themselves independent.

\end{frame}

%%%%%%%%%%%%%%%%%%%%%%%%%%%%%%%%%%%%%%%%%%%%%%%%%%%
\begin{frame}[fragile] \frametitle{}

{\bf F-test}

Now consider the hypothesis test $H_0: D\beta = d$ for a
matrix $D$ with $k$ columns and rank $k$.

\end{frame}

%%%%%%%%%%%%%%%%%%%%%%%%%%%%%%%%%%%%%%%%%%%%%%%%%%%
\begin{frame}[fragile] \frametitle{}

We'll form the following test statistic:
\begin{align*}
F &= \frac{(D\widehat{\beta} - d)^t [D (X^t X)^{-1} D^t ]^{-1} (D\widehat{\beta} - d) / k}{s^2}
\end{align*}
And prove that is has an F-distribution with $k$ and $n-p$ degrees of
freedom.

\end{frame}

%%%%%%%%%%%%%%%%%%%%%%%%%%%%%%%%%%%%%%%%%%%%%%%%%%%
\begin{frame}[fragile] \frametitle{}

We can re-write the test statistics as:
\begin{align*}
F &= \frac{w / k}{q / (n - p)}
\end{align*}
Where:
\begin{align*}
w &= (D\widehat{\beta} - d)^t [D (X^t X)^{-1} D^t ]^{-1} (D\widehat{\beta} - d) \\
q &= r^t r / \sigma^2
\end{align*}

\pause
We've already shown that $q | X \sim \chi^2_{n-p}$, and can see
that $q \independent w$ by the same argument as before. All that
is left is to show that $w \sim \chi^2_{k}$.

\end{frame}

%%%%%%%%%%%%%%%%%%%%%%%%%%%%%%%%%%%%%%%%%%%%%%%%%%%
\begin{frame}[fragile] \frametitle{}

Let $v = D\widehat{\beta} - d$. Under the null hypothesis $v = D (\widehat{\beta} - \beta)$.
Therefore:
\begin{eqnarray*}
\mathbb{V} (v | X) &=& \mathbb{V} (D (\widehat{\beta} - \beta)) \\ \pause
&=& D \mathbb{V} (\widehat{\beta} - \beta | X) D^t \\ \pause
&=& \sigma^2 D (X^t X)^{-1} D^t
\end{eqnarray*}

\pause
Now, for the multivariate normally distributed $v$ with zero mean,
we have:
\begin{align*}
w &= v^t \mathbb{V} (v | X)^{-1} v
\end{align*}
And therefore $w | X \sim \chi^2_k$

\end{frame}

%%%%%%%%%%%%%%%%%%%%%%%%%%%%%%%%%%%%%%%%%%%%%%%%%%%
\begin{frame}[fragile] \frametitle{}

{\bf An alternative F-test}

Now, consider the following estimator:
\begin{align*}
\widetilde{\beta} &= \argmin_b || y - Xb ||_2^2 , \quad \text{s.t.} \quad Db = d
\end{align*}
\pause We then define the restricted residuals
$\widetilde{r} = y - X\tilde{\beta}$.

\end{frame}

%%%%%%%%%%%%%%%%%%%%%%%%%%%%%%%%%%%%%%%%%%%%%%%%%%%
\begin{frame}[fragile] \frametitle{}

{\bf An alternative F-test}

An alternative expression for the $F$-test statistic is:
\begin{align*}
F &= \frac{(\tilde{r}^t \tilde{r} -  r^t r) / k }{r^t r / (n - p)}
\end{align*}
Conceptually, it should make sense that this is large whenever
the null hypothesis is false.

\end{frame}

\end{document}











